\section{Fuentes de datos, parámetros y supuestos}
\label{anx:datos}
% Origen: 03-paso2-datos-modelo-conceptual.tex (fuentes, parámetros y supuestos) y el anexo de días hábiles, movidos al anexo.

\subsection{Validaciones por molinete (SBASE)}
\label{anx:sbase}

La fuente cuantitativa principal es el dataset abierto \emph{Subte: Viajes Molinetes}
\parencite{sbase2026molinetes}. Tiene una fila por molinete y por ventana de 15 minutos, con la línea, la
estación, el identificador del molinete y los pasajeros por tipo de pago. El identificador distingue los
dos vestíbulos de Constitución:
\begin{center}
\small
\begin{tabular}{@{}ll@{}}
Vestíbulo Principal & \texttt{LineaC\_Constitucion\_Turn\textit{NN}} \\
Vestíbulo Plaza & \texttt{LineaC\_Constitucion\_Plaza\_Turn\textit{NN}} \\
\end{tabular}
\end{center}
Por eso la unidad del dato coincide con la unidad de servicio del modelo: el molinete.

El procesamiento, reproducible con el script \texttt{sbase-perfil.py} del repositorio, hizo cuatro cosas:
\begin{enumerate}
  \item \textbf{Normalizó el campo horario.} En marzo y abril de 2026 viene como \texttt{HH:MM} y en el
        resto de los meses como \texttt{HH:MM:SS}. Sin normalizar, cada hora se parte en dos grupos y los
        totales salen mal sin ningún aviso.
  \item \textbf{Excluyó los días atípicos.} Se descartaron los días cuyo flujo diario quedó por debajo del
        50~\% de la mediana. Los días excluidos coinciden uno a uno con feriados y fines de semana largos,
        lo que valida el criterio.
  \item \textbf{Fijó el período representativo.} Se usaron los días hábiles de marzo a junio de 2026:
        80 días útiles. Enero y febrero se descartaron porque son meses de vacaciones con menor carga.
  \item \textbf{Separó los vestíbulos.} El Principal tiene 22 identificadores: \texttt{Turn07} y de
        \texttt{Turn09} a \texttt{Turn29}. \texttt{Turn07} registró 3 pasajeros en seis meses, por lo que
        en la práctica hay 21 molinetes con tráfico. Por ventana se observan en promedio 19,5 molinetes
        activos.
\end{enumerate}

La figura~\ref{fig:perfil} muestra el perfil resultante junto con los arribos del Roca. Entre enero y junio,
\texttt{Turn14} procesó \num{1485} pasajeros por día hábil en la franja y \texttt{Turn23}, 129: un factor de 11.

\begin{figure}[H]
\centering
\includegraphics[width=0.85\textwidth]{tpi-subte-perfil-demanda.png}
\caption{Validaciones medias por ventana en el vestíbulo Principal (barras; en rojo, el pico) y arribos de
los trenes del Roca a Plaza Constitución según el horario oficial (marcas). Elaboración propia sobre
\textcite{sbase2026molinetes} y \textcite{trenes2026roca}.}
\label{fig:perfil}
\end{figure}

El dato tiene dos límites, y ambos condicionan el modelo:
\begin{itemize}
  \item \textbf{Cuenta validaciones, no arribos ni esperas.} No dice cuánta gente esperó ni cuánto
        tiempo. Si en algún momento los molinetes se saturan, el conteo mide la capacidad de los
        molinetes y no la demanda real (censura por capacidad).
  \item \textbf{Es un agregado de 15 minutos.} Dentro de cada ventana se borra la estructura de las
        oleadas, que es justamente lo que genera la cola.
\end{itemize}

\subsection{Horario del Roca}

La estructura temporal de las oleadas se obtuvo del horario oficial de la línea Roca, vigente desde el
03/08/2026 \parencite{trenes2026roca}. Se procesaron los cuatro cuadros horarios que llegan a Plaza
Constitución: La Plata, Bosques, Glew/A.~Korn y Ezeiza/Cañuelas. De cada uno se tomaron la sección de
lunes a viernes, los trenes con destino Constitución y el horario de llegada. Se descartaron las
lanzaderas que no llegan a la terminal y se controló que un tren repetido en dos cuadros tuviera el mismo
arribo. Quedan 47 trenes entre las 07:00 y las 09:30, uno cada 3 minutos en mediana, con huecos de hasta
8 minutos. Cuatro pares de trenes llegan en el mismo minuto.

El horario fija cuándo llega cada oleada, pero no cuántas personas trae ni qué parte de ellas va a la
Línea~C.

\subsection{Pedidos de información pública}
\label{anx:pedidos}

Lo que ningún dataset publica se pidió a los organismos responsables. El 23/09/2026 se enviaron correos
a SBASE, a Emova (la concesionaria) y a Trenes Argentinos. El 05/10 se presentaron pedidos formales:
\begin{itemize}
  \item a SBASE, por la Ley~104 de la Ciudad: el plano de los vestíbulos, la correspondencia entre
        identificadores y equipos, el criterio de habilitación de molinetes y los tiempos de validación;
  \item a Trenes Argentinos Operaciones, por la Ley~\num{27275}: el plano de la estación, los pasajeros
        por tren, las llegadas efectivas y los estudios de trasbordo.
\end{itemize}

Trenes Argentinos respondió el 26/10/2026 \parencite{sofse2026nota}. Envió:
\begin{itemize}
  \item el plano de la estación ferroviaria en DWG, que no incluye los accesos de la Línea~C;
  \item la llegada real de cada tren entre el 02/03 y el 30/06/2026, de 07:00 a 09:30;
  \item una matriz origen-destino de la línea Roca para el mismo período.
\end{itemize}
Informó además que no existen estudios de trasbordo hacia el subte. Estos datos llegaron después de la
primera producción, la del 05/10 con un día promedio. El modelo final los usa como entradas: el plano fija los
andenes, el hall y las escaleras mecánicas, y cada réplica toma las llegadas reales de su día. La matriz
origen-destino se usa en la validación (anexo~\ref{anx:validacion}). SBASE prorrogó su plazo de
respuesta y, al cierre de este informe, no había enviado el plano del subsuelo: la plaza de transbordo y el
interior del vestíbulo siguen siendo supuestos.

\subsection{Parámetros sin dato publicado: fuentes}
\label{anx:parametros-fuentes}

La tabla~\ref{tab:parametros-fuentes} detalla la fuente y el carácter de cada rango de la
tabla~\ref{tab:parametros}.

\begin{table}[H]
\centering
\caption{Parámetros sin dato publicado: fuente y carácter de cada rango}
\label{tab:parametros-fuentes}
\small
\begin{tabularx}{\textwidth}{@{}>{\raggedright\arraybackslash}p{3.0cm} X@{}}
\toprule
Parámetro & Fuente y carácter \\
\midrule
Tiempo de validación &
Moda: 2,4~s con tarjeta inteligente en Londres; 2,6 a 2,9~s en Nueva York \parencite[Exh.~10-27]{tcqsm2013}.
Cota del propio dataset: hasta 246 pasajeros en 15 minutos por molinete, o sea una media sostenida
$\le 900/246 = 3{,}66$~s. Asimetría a la derecha por fallas de lectura [supuesto] \\
Descenso por las puertas &
Tiempo en que el pasaje de un tren sale por las 21 puertas del lado del andén (7 coches con 3 puertas
por lado, \cite{wikicsr}): con 200 a 530 pasajeros por tren, unos 25 por puerta a 1 a 1,5 por segundo
[supuesto]. El andén ya no entra acá: se camina. Una descarga corta concentra la oleada y genera más cola \\
Demora hasta abrir las puertas &
Desde el arribo del horario hasta que se abren las puertas [supuesto]. Antes era la demora andén--vestíbulo,
$D/v + h/v_e$ con $D = 85$~m (\cite{osm2026}); la superposición con el plano llevó a simular el andén, el
hall, las escaleras al subsuelo y el pasillo, y la demora quedó solo para la apertura (anexo~\ref{anx:geometria}) \\
Proporción que llega en tren &
Supuesto explícito: no hay estudio publicado de trasbordo del Roca a la Línea~C, y Trenes Argentinos
confirmó que no lo tiene \parencite{sofse2026nota}. Se intentó estimarla cruzando el dato de SBASE con
las llegadas de trenes por ventana; no es identificable (anexo~\ref{anx:volumen}) \\
\bottomrule
\end{tabularx}
\end{table}

La demora solo desplaza las oleadas en el tiempo, sin cambiar su forma, por lo que se fija en el valor
intermedio. Los otros tres parámetros moldean la congestión y son los candidatos al análisis de
sensibilidad.

La velocidad de marcha se toma uniforme entre 1,1 y 1,5~m/s, el rango que \textcite{li2024} reportan en
la literatura y dentro del cual cae su valor calibrado de 1,37~m/s. El diámetro del peatón se fijó en 0,5~m.
Es un supuesto: ninguno de los estudios citados lo informa.

\subsection{Generación de la demanda}
\label{anx:demanda}

La demanda de cada ventana $k$ es la observada $d_k$ del día de la
réplica, y define una tasa de llegada a lo largo de la franja. Una fracción $p$ (la proporción que llega en tren) viaja en los
trenes del Roca. La gente espera en las estaciones de su ramal, así que cada tren lleva la demanda que se
acumuló desde el tren anterior \emph{de su mismo ramal}. Cada ramal toma la parte de la demanda proporcional
a sus trenes en la franja. Con este criterio, un tren que llega después de un intervalo largo viene más
cargado, y dos trenes de ramales distintos que llegan juntos no se reparten la misma gente. Esa demanda
acumulada es la media del pasaje del tren: la cantidad efectiva se sortea con una distribución de Poisson
[supuesto], porque la carga de cada formación varía aunque la tasa sea la misma. El total de la
franja coincide con el dato; ventana por ventana, el ingreso se corre hacia el tren siguiente.

El pasaje de cada tren baja por las 21 puertas de su lado del andén, de forma pareja durante el descenso,
y desde ahí camina el andén y el hall, baja por la escalera del bloque central que le queda de su lado y sigue
por la plaza de transbordo y el pasillo hasta el vestíbulo. El resto, $(1-p)\,d_k$, entra desde la
calle con llegadas de Poisson no homogéneas a la tasa de cada ventana. Solo se generan ingresos antes de las
09:30.

\subsection{Supuestos del modelo}
\label{anx:supuestos}

Además de los rangos de la tabla~\ref{tab:parametros}, el modelo adopta estos supuestos:
\begin{enumerate}
  \item La cola está vacía a las 07:00.
  \item Los molinetes funcionan solo en sentido de ingreso durante la franja, y los habilitados están
        disponibles todo el período.
  \item Nadie abandona la cola ni cambia de molinete una vez elegido.
  \item El volumen de cada ventana y los trenes son los de un día hábil real con servicio normal; los días
        con el servicio del Roca interrumpido quedan fuera del estudio.
  \item El pasaje de cada tren es una variable de Poisson cuya media es la demanda acumulada de su ramal.
  \item El hall, el andén y las escaleras del bloque central siguen el plano de Trenes Argentinos. La plaza
        de transbordo, el ancho del pasillo y el interior del vestíbulo son supuestos
        (anexo~\ref{anx:recorrido}), porque SBASE no envió el plano del subsuelo.
  \item Todo pasajero del Roca que entra al modelo baja por las escaleras mecánicas; las escaleras fijas del
        bloque central no se modelan.
\end{enumerate}
El supuesto~4 reemplaza a dos versiones anteriores (perfil medio en todas las réplicas y trenes en
horario), que la validación mostró optimistas: un atraso variable junta trenes que el horario separa.

\subsection{Días hábiles y trenes del Roca}
\label{anx:roca}

Los 75 días del estudio son los días hábiles de marzo a junio de 2026 con datos de SBASE y servicio normal del
Roca. Se excluyeron cinco días con menos de 40 trenes en la franja, por servicio interrumpido: 23/03, 07/04,
13/04, 26/05 y 18/06. Las llegadas de cada tren a Plaza Constitución provienen de la respuesta de Trenes
Argentinos \parencite{sofse2026nota}; los perfiles por ventana, del dataset de SBASE. Los dos se embeben en el
modelo con \texttt{datos/embeber\_dias\_sbase.py}.
